%Simple_math_chapter.tex
%all the math used by BLESS from first principles

% !TEX root=../BLESS_Language_Reference_Manual.tex
%
%\pn This chapter provides an informal description of the mathematical concepts used by BLESS.
%These informal descriptions should be enough for engineers and programmers to understand the mathematics used by BLESS.
%For those who want formality and rigor for a deeper understanding, and to assure themselves that BLESS has a sound foundation, 
%chapters \ref{chapter:math} and \ref{chapter:BL} provide formality and rigor.
%
%\pn To prove correctness, a programming language must be mathematically defined.  Therefore, the foundational mathematics must be derived from First Principles.
%%\footnote{Much of the following was adapted from U.S. Pat. No. 5,867,649 DANCE/Multitude Concurrent Computation}
%
%\pn The foundational mathematics was deliberately selected to be as simple as possible, using only a few fundamental concepts from which all else flows.  The following sections tersely declare these fundamental concepts, and is not meant as a textbook or tutorial.   Many pieces of standard mathematics, like numbers and arithmetic are assumed as defined by \cite{CRCmath}.
%%\footnote{as defined by \emph{CRC Concise Encyclopedia of Mathematics}, Eric W Weisstein, editor, second edition, Chapman \& Hall/CRC, 2003.}
%
%\pn Later, computation will be defines as satisfaction of interval-temporal logic formulas with lattices of states--not as a sequence of imperative commands.  A lattice is a graph with some special properties.  Thinking about programs as logic formulas, instead of traditional, imperative, sequential control flow, takes some getting used to.
%
%\pn Still, this document attempts to be self-contained, explicitly built on simple math defined herein, starting with sets.
 

\section{Sets}\label{sec:informal_sets}

\pn A \bemph{set} is a collection of elements.  Finite sets may be specified by enumerating their elements between curly braces.  For example, \(\{true, false\}\) denotes the set consisting of the Boolean constants \bemph{true} and \bemph{false}.  When enumerating elements of a set, ``\ldots" is used to denote repetition.  For example, \(\{1, 
\ldots, n\}\) denotes the set of natural numbers from 1 to $n$ where the upper bound, $n$, is a natural number that is not further specified.

\pn More generally, classes (which are almost always sets, too) are specified by referring to some property of their elements.
\(\{x~\vert~ P\}\)
denotes the set of all elements $x$ that satisfy the property $P$.  The bar, $~\vert~$
\index{ bar@$~\vert~$ such that}, can be read as ``such that".
For example,
\(\{x~\vert~x~$is an integer and$~x~$is divisible by 2$\}\)
denotes the infinite set of all even integers.

\pn For membership, $a\in A$\index{ in@$\in$ element of set} denotes that $a$ is an element of the set $A$, and 
$b\notin A$\index{ notin@$\notin$ not element of set} to denote that $b$ is not an element of the set $A$.  

\pn Some sets have customary symbols: 
\begin{description}
\item[$\emptyset$]\index{ emptyset@$\emptyset$ empty set} denotes the empty 
set;
%\footnote{Some claim that $\emptyset$ is not unique.  The set of sports cars I own, and the set of supermodels I've dated are both empty, but are different sets.  Others deny the existence of $\emptyset$ claiming that a set that is empty is nothing at all.} 
\item[$\mathbb{N}$]\index{ natural@$\mathbb{N}$ natural number} denotes the set of all natural numbers, excluding 0; 
\item[$\mathbb{N}_0$]\index{ natural@$\mathbb{N}_0$ natural} denotes the set of all natural numbers, including 0; 
\item[$\mathbb{Z}$]\index{ integer@$\mathbb{Z}$ integer} denotes the set of all integers; 
\item[$\mathbb{Q}$]\index{ rational@$\mathbb{Q}$ rational} denotes the set of rational numbers; 
\item[$\mathbb{R}$]\index{ real@$\mathbb{R}$ real} denotes the set of real numbers;
\item[$\mathbb{C}$]\index{ complex@$\mathbb{C}$ complex} denotes the set of complex numbers.
\item[$\mathbb{B}$]\index{ boolean@$\mathbb{B}$ boolean} denotes the set \(\{\top, \bot\}\) \(\{{true}, {false}\}\).
\end{description}
Fixed-point\index{fixed-point} numbers are rational numbers with fixed divisor.

\pn In a set, one does not distinguish repetitions of elements.  Thus \(\{T,F\}\) and \(\{T,T,F\}\) are the same set.
Often it is convenient to refer to a given set (or class) when defining a new set (or class).
\(\{x\in A~\vert~ \psi\}\) (where $\psi$ is some well-formed formula containing $x$, and $\in$ means member-of) is an abbreviation for \(\{x~\vert~ x\in A~$and$~\psi\}\).  Similarly, the order of elements is irrelevant.  Two sets $A$ and $B$ are equal, $B=A$, if-and-only-if they have the same elements.

\pn Let $A$ and $B$ be sets.  Then $A\subseteq B$\index{ subset@$\subseteq$ subset} denotes that $A$ is a \bemph{subset} of $B$; $A\cap B$\index{ intersection@$\cap$ intersection} denotes the \bemph{intersection} of $A$ and $B$; 
$A\cup B$\index{ union@$\cup$ union} denotes the \bemph{union} of $A$ and $B$; and, $A \minus B$ denotes the \bemph{difference} of $A$ and $B$.  The symbol $\equiv$\index{ equivalence@$\equiv$ equivalence} is used to define equivalence.

\(A \subseteq B \equiv a \in B~$for every$~a \in A\) \\
\(A\cap B \equiv \{a~\vert~ a\in A~$and$~a\in B\}\) \\
\(A\cup B \equiv \{a~\vert~ a\in A~$or$~a\in B\}\) \\
\(A \minus B \equiv \{a~\vert~ a\in A~$and$~a\notin B\}\) 

\pn Sets $A$ and $B$ are \bemph{disjoint} if they have no element in common, $A\cap B=\emptyset$, where $\emptyset$ is the empty set..

\pn The definitions of intersection and union can be generalized to more than two sets.  Let $A_k$ be a set for every element $k$ of some other set $J$ in
\(\bigcap_{k\in J} \equiv \{a~\vert~ a\in A_k~$for all$~k\in J\}\)
\(\bigcup_{k\in J} \equiv \{a~\vert~ a\in A_k~$for some$~k\in J\}\)

\pn For a finite set $A$, $\Vert A\Vert$\index{ cardinality@$\Vert A\Vert$ cardinality} denotes the \bemph{cardinality}, or number of elements in $A$.\footnote{KerML uses \# as a cardinality operator.}  For a non-empty, finite set $B\subset\mathbb{Z}$, $\bmin(B)$ denotes the \bemph{minimum} of all integers in $B$.

\pn $\mathbb{D}$ is the set of all possible constructed values, including strings, records, and arrays.% formally defined in Chapter \ref{chapter:type}.

\section{Tuples}\label{sec:informal_tuples}

\pn For sets, the repetition of elements and their order is irrelevant.  When ordering matters, ordered pairs and tuples are used.  For elements $a$ and $b$, not necessarily distinct, $\langle a,b\rangle $ is an \bemph{ordered pair} or simply pair.  Then $a$ and $b$ are called \bemph{components} of $\langle a,b\rangle$.  Two pairs $\langle a,b\rangle$ and $\langle b,c\rangle$ are equal,
$\langle a,b\rangle=\langle c,d\rangle$ if-and-only-if $a=c$ and $b=d$.

\pn More generally, let $n$ be any natural number (excluding 0), $n\in\mathbb{N}$.  Then if $a_1,\ldots,a_n$ are any $n$ elements, then $\langle a_1,\ldots,a_n\rangle$ is an \bemph{n-tuple}.   The element $a_k$ where $k\in\{1,\ldots,n\}$ is called the k-th element of $\langle a_1,\ldots,a_n\rangle$.
An n-tuple $\langle a_1,\ldots,a_n\rangle$ is equal to an m-tuple 
$\langle b_1,\ldots,b_m\rangle$, $\langle a_1,\ldots,a_n\rangle=\langle b_1,\ldots,b_m\rangle$, if-and-only-if $m=n$ and $a_k=b_k$ for all $k\in\{1,\ldots,n\}$.  Note that 2-tuples are pairs.  Additionally, a 0-tuple is written as $\langle \rangle$, and a 1-tuple as $\langle a\rangle$ for any element $a$.

\pn The \bemph{Cartesian product}, $A\times B$\index{ product@$\times$ product} of sets $A$ and $B$ consists of all pairs, $\langle a,b\rangle$ with $a\in A$ and $b\in B$.  The n-fold Cartesian product, $A_1\times\ldots\times A_n$ of sets $A_1,\ldots,A_n$ consists of all n-tuples, 
$\langle a_1,\ldots,a_n\rangle$ with $a_k\in A_k$ for $k\in\{1,\ldots,n\}$.  If all $A_k$ are the same set $A$, then the n-fold Cartesian product, $A\times\ldots\times A$ is also written $A^n$.

\section{Relations}\label{sec:informal_relationsmath}

\pn A binary \bemph{relation} $R$ between sets $A$ and $B$ is a subset of their Cartesian product, $A\times B$, that is, 
$R\subseteq A\times B$.  If $A=B$, then $R$ is called a relation on $A$.  For example,
\(\{\langle a,1\rangle,\langle b,2\rangle,\langle c,2\rangle\}\)
is a binary relation between $\{a,b,c\}$ and $\{1,2\}$.  More generally, for any natural number $n$, and n-ary relation $R$ between sets $A_1,\ldots,A_n$ is a subset of the n-fold Cartesian product $A_1\times\ldots\times A_n$, that is, 
$R\subseteq A_1\times\ldots\times A_n$.  Note that 1-ary relations are called unary relations, 2-ary relations are called binary relations, and 3-ary relations are called ternary relations.

\pn Consider a binary relation $R$ on a set $A$.  $R$ is called \bemph{reflexive} if $\langle a,a\rangle\in R$ for every $a\in A$; it is called \bemph{irreflexive} if $\langle a,a\rangle\notin R$ for every $a\in A$.  $R$ is called \bemph{symmetric} if for all $a,b\in A$, whenever $\langle a,b\rangle\in R$ the also $\langle b,a\rangle\in R$; it is called \bemph{antisymmetric} if for all $a,b\in A$, whenever 
$\langle a,b\rangle\in R$ and $\langle b,a\rangle\in R$ then $b=a$.  $R$ is called \bemph{transitive} if for all $a,b,c\in A$ whenever $\langle a,b\rangle\in R$ and $\langle b,c\rangle\in R$ then also $\langle a,c\rangle\in R$.

\pn The transitive, reflexive \bemph{closure}, $R^\star$\index{ closure@$R^{\star}$ transitive reflexive closure}, of a binary relation $R$ over a set $A$, is the smallest, transitive and reflexive, binary relation on $A$ that contains $R$ as a subset.  The transitive, irreflexive closure, $R^{+}$\index{ closure2@$R^{+}$ transitive irreflexive closure}, 
 of a binary relation $R$ over a set $A$, is the smallest, transitive and irreflexive binary relation that contains $R$ as a subset.
 
\(R^\star \equiv~R \subseteq R^\star~\band~$for all$~a,b,c\in A ~\vert~  
 \begin{array}{c}
 \langle a,b\rangle\in R^\star\wedge\langle b,c\rangle\in R^\star \rightarrow \langle a,c\rangle\in R^\star\\
 \langle a,a\rangle\in R^\star
 \end{array}
 \)
 
\(R^{+} \equiv~R \subseteq R^{+}~\band~$for all$~a,b,c\in A ~\vert~  
 \begin{array}{c}
 \langle a,b\rangle\in R^{+}\wedge\langle b,c\rangle\in R^{+} \rightarrow \langle a,c\rangle\in R^{+}\\
 \langle a,a\rangle\notin R^{+}
 \end{array}
 \)

\pn The \bemph{relational composition}, $R_1\circ R_2$
\index{ composition@$\circ$ relational composition}, of relations $R_1$ and $R_2$ on a set $A$ creates a new relation by combining them:

\(R_1\circ R_2 \equiv \{\langle a,c\rangle~\vert~~$there exists$~b\in A~$with$~\langle a,b\rangle\in R_1 \band~\langle b,c\rangle\in R_2\}\)

\pn For any natural number $n$, the n-fold relational composition, $R^{n}$, of a relation $R$ on a set $A$ is defined inductively:

\(R^{0} \equiv \{\langle a,a\rangle~\vert~ a\in A\}\)

\(R^{n} \equiv R^{n\minus 1}\circ R\)
for $n>0$.  

\(R^\star \equiv \bigcup_{n\in \mathbb{N}_0} R^{n}\)

\(R^{+} \equiv R^\star \minus R^{0}\)

\pn Membership of pairs in a binary relation is usually written in infix notation; instead of $\langle a,b\rangle\in R$, use $aRb$.
Any binary relation $R\subseteq A\times B$ has an inverse $R^{\minus1} \subseteq B\times A$ such that 
$b R^{\minus1}a$ if-and-only-if $a R b$.

\section{Functions}\label{sec:informal_functions}

\pn Let $A$ and $B$ be sets.  A \bemph{function} or mapping from $A$ to $B$ is a binary relation $f$
between $A$ and $B$ with the following special property:  for each element $a\in A$ there is exactly one element $b\in B$ such that $afb$.  Usually functions use prefix notation for function application writing $f(a)=b$ 
instead of $afb$.  For some functions postfix notation is used to write $af=b$.  To indicate that $f$ is a 
function from $A$ to $B$ write 
\(f:A\rightarrow B\).    
The set $A$ is called the \bemph{domain} of $f$ and the set $B$ is called the \bemph{range} or \bemph{co-domain} of $f$.

\pn Consider a function \(f:A\rightarrow B\) and some set $X\subseteq A$.  The \bemph{restriction} of $f$ to $X$ 
is denoted by $f[X]$ and defined as the intersection of $f$ (which is a subset of $A\times B$) with $X\times B$:
\(f[X] \equiv f\cap(X\times B)\).  
Functions may have special properties.  A function \(f:A\rightarrow B\) is called \bemph{one-to-one} or \bemph{injective} if $f(a_1)\ne f(a_2)$ for any two distinct elements $a_1,a_2\in A$.  It is called \bemph{onto} or 
\bemph{subjective} if for every element $b\in B$ there exists an element $a\in A$ with $f(a)=b$.  It is called 
\bemph{bijective} if it is both injective and subjective.

\pn Consider an n-ary function whose domain is a Cartesian product, $f:A_1 \times\ldots\times A_n~\rightarrow~B$.  It is customary to drop tuple brackets when applying $f$ to a tuple\index{tuple} 
\(\langle a_1,\ldots,a_n\rangle\in~A_1 \times\ldots\times A_n\) 
writing $f(2,3)$ instead of $f(\langle 2,3\rangle)$.

\pn Consider a binary function whose domain and co-domain coincide, $f:A\rightarrow A$.  An element
$a\in A$ is called a \bemph{fixed point} of $f$ if $f(a)=a$.


\pn Boolean logic can also be considered to be functions on \{\emph{true}, \emph{false}\}.
\begin{description}
\item[\bemph{conjunction}] of $a$ and $b$ is $a\wedge b$\index{ conjunction@$\wedge$ conjunction}; 
\item[\bemph{disjunction}] is $a\vee b$ \index{ disjunction@$\vee$ disjunction}; 
\item[\bemph{implication}] is $a\rightarrow b$\index{ implication@$\rightarrow$ implication}; 
\item[\bemph{if-and-only-if}] is $a\leftrightarrow b$ \index{ iff@$\leftrightarrow$ if-and-only-if}; 
\item[\bemph{exclusive disjunction}] is $a\oplus b$\index{ exclusive@$\oplus$ exclusive disjunction};
\item[\bemph{complement}] is $\neg a$\index{ complement@$\neg$ complement}.
\end{description}

The meaning of boolean functions is defined by truth table \ref{booleantt}.

\begin{table}[htp]
\caption{Boolean Function Truth Table}
\begin{center}
\begin{tabular}{|c|c||c|c|c|c|c|c|}
\hline
$a$ & $b$ & $a\wedge b$ & $a\vee b$ & $a\rightarrow b$ & $a\leftrightarrow b$ & $a\oplus b$ & $\neg a$ \\  \hline 
\emph{false} & \emph{false} & \emph{false} & \emph{false} & \emph{true} & \emph{true} & \emph{false} &  \emph{true}  \\ \hline 
\emph{true} & \emph{false} & \emph{false} & \emph{true} & \emph{false} & \emph{false} & \emph{true} & \emph{false}   \\ \hline 
\emph{false} & \emph{true} & \emph{false} & \emph{true} & \emph{true} & \emph{false} & \emph{true} &  \emph{true}  \\ \hline 
\emph{true} & \emph{true} & \emph{true} & \emph{true} & \emph{true} & \emph{true} & \emph{false} &  \emph{false}  \\ \hline 
\end{tabular}
\end{center}
\label{booleantt}
\end{table}%

\section{Sequences}\label{sec:informal_sequences}

\pn Sequences are like tuples in that their elements are ordered, but sequences are not constrained to have a specific number of elements.

\pn In the following, let $A$ be a set.  A \bemph{sequence} of elements of $A$ of length $n>0$ is a function
\(f:\{1,\ldots,n\} \rightarrow A\).  
A sequence is denoted by listing its values in order
\(a_1,\ldots,a_n\)
where $a_1=f(1),~\ldots~, a_n=f(n)$.  Then the k-th element of the sequence $a_1,\ldots,a_n$ is $a_k$ when 
$k\in\{1,\ldots,n\}$.  A finite sequence is a sequence of any length $n\ge0$.
A sequence of length 0 is called the \bemph{empty sequence} and denoted $\epsilon$\index{ epsilon@$\epsilon$ empty sequence}.
A countably-infinite sequence of elements from A is a function $\xi:\mathbb{N}_0~\rightarrow~A$.
To exhibit the general form of a countably-infinite sequence $\xi$ is written
\(\xi:a_0~a_1~a_2~\ldots\)
when $a_k=\xi(k)$ for all $k\in\mathbb{N}_0$.  Then $k$ is called the \bemph{index} of element $a_k$.

\pn Consider now a set of binary relations, $R=\{R_1,R_2,\ldots,R_{n\minus 1}\}$, on a set $A$.  For any finite sequence of elements of $A$, 
$a_1~\ldots~a_n$, such that each element is related to the next by a relation in $R$,
\(a_1 R_1 a_2,~a_2 R_2 a_3,~\ldots,~a_{n\minus 1} R_{n\minus 1} a_n\)
can be written as a finite chain,
\(a_1 R_1 a_2 R_2 a_3 \ldots R_{n\minus 1} a_n\)
For example, using the relations $=$ and $<$ over $\mathbb{Z}$, a finite chain may be written
\(a_0<a_1=a_2<a_3<a_4\).  
Similarly for infinite sequences and infinite chains.

\pn A \bemph{permutation} of a sequence has the same elements in different order.  In the following, let $f$ and $g$ be sequences of distinct\footnote{may not need restriction on repeated elements; that the sets are the same, repeated elements and all, may be enough; but then they have equal bags, not sets and that way madness lies} elements
\(f:\{1,\ldots,n\} \rightarrow A\) and \(g:\{1,\ldots,m\} \rightarrow A\).  
The sequences are permutations of each other, $f \leftrightarrows g$,\index{ permutation@$\leftrightarrows$ permutation} when they are the same length and have the same elements
\(f\leftrightarrows g~\equiv~n=m\land \{f(1),f(2),\ldots,f(n)\}=\{g(1),g(2),\ldots,g(m)\}\)


\section{Strings}\label{sec:informal_strings}

\pn A set of symbols is often called an \bemph{alphabet}.  A \bemph{string} over an alphabet $A$
is a finite sequence of symbols from $A$.  For example, $1+2$ is a string over the alphabet
$\{1,2,+\}$.  %Syntactic objects like AADL annex subclauses are strings.

\pn The \bemph{concatenation} of two strings $s_1$ and $s_2$ yields the string $s_{1}s_{2}$ 
formed by first writing $s_1$ and then $s_2$.  
A string $t$ is called a \bemph{substring} of a string $s$ if there exist strings $s_1$ and $s_2$
such that $s=s_{1}t s_{2}$.  Because $s_1$ and $s_2$ may be empty, every string is a substring of itself.

\section{Partial Orders}\label{sec:informal_partialorders}

\pn A \bemph{partial order} is a pair $\langle A, \sqsubset\rangle$\index{ partial@$\sqsubset$} consisting of a set $A$ and a irreflexive, antisymmetric, and transitive relation $\sqsubset$ on $A$. The reflexive partial order is denoted $\sqsubseteq$\index{ partialeq@$\sqsubseteq$} .  
If $x\sqsubset y$ for some $x,y\in A$, then $x$ is called \bemph{less than} $y$, or $y$ is \bemph{greater than} $x$.  
Consider an element $a\in A$ and a subset $X\subseteq A$.  When $a\in X$ and $a \sqsubset x$ for all $x\in X\minus\{a\}$, the $a$ is called the \bemph{least element} of $X$.  
When $x \sqsubset b$ for all $x\in X\minus\{b\}$, then $b$ is called an \bemph{upper bound} of 
$X$.  Upper bounds of $X$ need not be elements of $X$.
Let $U$ be the set of all upper bounds of $X$.  Then $a$ is called the \bemph{least upper bound} of $X$ if 
$a$ is the least element of $U$.

\pn A partial order $(A, \sqsubset)$ is called \bemph{complete} if $A$ contains a least element, and 
for every ascending chain
\(a_0 \sqsubset a_1 \sqsubset a_2 \cdots\)
of elements from A, the set
\(\{a_0,a_1,a_2,\ldots\}\)
has a least upper bound.

\section{Graphs}\label{sec:informal_graphs}

\pn A \bemph{graph} is a pair $\langle V,E\rangle$ where $V$ is a finite set of vertices
\(\{v_1,v_2,\ldots,v_n\}\)
and $E$ is a finite set of edges where each edge is a pair of vertices in $V$,
\(\{\langle v_m,v_l\rangle,\ldots,\langle v_j,v_k\rangle\}\).  
All graphs considered here are \bemph{directed} in the the order of vertices within the pair
describing the edge is significant.  
The set of edges forms a relation on the set of vertices
\(E\subseteq V\times V\).  
The transitive, irreflexive closure of $E$, called $E^{+}$ is especially important.

%\section{Lattices}\label{sec:informal_lattices}
%
%\pn A graph $\langle V,E\rangle$ is a \bemph{lattice} if the transitive, irreflexive closure of $E$, $E^{+}$,
%is an irreflexive partial order $\sqsubset$\index{ partial@$\sqsubset$}, it has a least element $\ell\in V$, and an upper
%bound $u\in V$.  Executions of \lang~ actions create lattices. 
%
%\pn Depictions of lattices place the least element (a.k.a. ``start") at the top, and the greatest element (a.k.a. ``end") at the bottom.
%Directed edges use no arrowheads; instead, the edges are presumed to flow from the higher vertex to the lower vertex.
%
%\setlength{\unitlength}{0.6pt}
%\renewcommand{\nodetextscalefactor}{1} %scales text with unitlength	
%
%%\begin{figure}
%\begin{wrapfigure}{l}{0pt}
%%lattice1.tex
%draws a pretty, generic lattice

\setcounter{lwidth}{390}
\setcounter{lheight}{300}

\begin{picture}(\value{lwidth},\value{lheight})(0,0)
%\graphpaper(0,0)(\value{lwidth},\value{lheight})

\setcounter{circlediameter}{12}
\thicklines
\color{blue}
% \drawarc{row}{numrows}{fromnode}{nodesinfromrow}{tonode}{nodesintorow}
\drawarc{1}{11}{1}{1}{1}{6}
\drawarc{1}{11}{1}{1}{2}{6}
\drawarc{1}{11}{1}{1}{3}{6}
\drawarc{1}{11}{1}{1}{4}{6}
\drawarc{1}{11}{1}{1}{5}{6}
\drawarc{1}{11}{1}{1}{6}{6}

\drawarc{2}{11}{1}{6}{1}{1}
\drawarc{2}{11}{2}{6}{1}{1}
\drawarc{2}{11}{3}{6}{1}{1}
\drawarc{2}{11}{4}{6}{1}{1}
\drawarc{2}{11}{5}{6}{1}{1}
\drawarc{2}{11}{6}{6}{1}{1}

\drawarc{3}{11}{1}{1}{1}{3}
\drawarc{3}{11}{1}{1}{2}{3}
\drawarc{3}{11}{1}{1}{3}{3}

\drawarc{4}{11}{1}{3}{1}{6}
\drawarc{4}{11}{1}{3}{2}{6}
\drawarc{4}{11}{2}{3}{3}{6}
\drawarc{4}{11}{2}{3}{4}{6}
\drawarc{4}{11}{3}{3}{5}{6}
\drawarc{4}{11}{3}{3}{6}{6}

\drawarc{5}{11}{1}{6}{1}{12}
\drawarc{5}{11}{1}{6}{2}{12}
\drawarc{5}{11}{1}{6}{3}{12}
\drawarc{5}{11}{2}{6}{4}{12}
\drawarc{5}{11}{3}{6}{5}{12}
\drawarc{5}{11}{3}{6}{6}{12}
\drawarc{5}{11}{3}{6}{7}{12}
\drawarc{5}{11}{4}{6}{8}{12}
\drawarc{5}{11}{5}{6}{9}{12}
\drawarc{5}{11}{5}{6}{10}{12}
\drawarc{5}{11}{5}{6}{11}{12}
\drawarc{5}{11}{6}{6}{12}{12}

\drawarc{6}{11}{1}{12}{1}{6}
\drawarc{6}{11}{2}{12}{1}{6}
\drawarc{6}{11}{3}{12}{1}{6}
\drawarc{6}{11}{4}{12}{2}{6}
\drawarc{6}{11}{5}{12}{3}{6}
\drawarc{6}{11}{6}{12}{3}{6}
\drawarc{6}{11}{7}{12}{3}{6}
\drawarc{6}{11}{8}{12}{4}{6}
\drawarc{6}{11}{9}{12}{5}{6}
\drawarc{6}{11}{10}{12}{5}{6}
\drawarc{6}{11}{11}{12}{5}{6}
\drawarc{6}{11}{12}{12}{6}{6}

\drawarc{7}{11}{1}{6}{1}{3}
\drawarc{7}{11}{2}{6}{1}{3}
\drawarc{7}{11}{3}{6}{2}{3}
\drawarc{7}{11}{4}{6}{2}{3}
\drawarc{7}{11}{5}{6}{3}{3}
\drawarc{7}{11}{6}{6}{3}{3}

\drawarc{8}{11}{1}{3}{1}{1}
\drawarc{8}{11}{2}{3}{1}{1}
\drawarc{8}{11}{3}{3}{1}{1}

\drawarc{9}{11}{1}{1}{1}{10}
\drawarc{9}{11}{1}{1}{2}{10}
\drawarc{9}{11}{1}{1}{3}{10}
\drawarc{9}{11}{1}{1}{4}{10}
\drawarc{9}{11}{1}{1}{5}{10}
\drawarc{9}{11}{1}{1}{6}{10}
\drawarc{9}{11}{1}{1}{7}{10}
\drawarc{9}{11}{1}{1}{8}{10}
\drawarc{9}{11}{1}{1}{9}{10}
\drawarc{9}{11}{1}{1}{10}{10}

\drawarc{10}{11}{1}{10}{1}{1}
\drawarc{10}{11}{2}{10}{1}{1}
\drawarc{10}{11}{3}{10}{1}{1}
\drawarc{10}{11}{4}{10}{1}{1}
\drawarc{10}{11}{5}{10}{1}{1}
\drawarc{10}{11}{6}{10}{1}{1}
\drawarc{10}{11}{7}{10}{1}{1}
\drawarc{10}{11}{8}{10}{1}{1}
\drawarc{10}{11}{9}{10}{1}{1}
\drawarc{10}{11}{10}{10}{1}{1}

\color{green}
%\drawnoderow{1}{11}{1}	%draw start node
%\drawtextnode{row}{numrows}{node}{numnodesinrow}{text}
\drawtextnode{1}{11}{1}{1}{}%{$\ell$}
\drawnoderow{2}{11}{6}
\drawnoderow{3}{11}{1}
\drawnoderow{4}{11}{3}
\drawnoderow{5}{11}{6}
\drawnoderow{6}{11}{12}
\drawnoderow{7}{11}{6}
\drawnoderow{8}{11}{3}
\drawnoderow{9}{11}{1}
\drawnoderow{10}{11}{10}
%\drawnoderow{11}{11}{1}
% \drawtextnode{row}{numrows}{node}{numnodesinrow}{text}
\drawtextnode	{11}{11}{1}{1}{}%{$u$}

%annotations
\color{black}
\put(215,285){$l=$start(\interval{i})}		%label "start(i)"
\put(214,285){\vector(-4,-1){22}}		%arrow to top node of lattice
\put(214,10){$u=$end(\interval{i})}		%label "end(i)"
\put(214,15){\vector(-4,1){22}}	%arrow to top node of lattice

\thinlines
\put(342,85){\scalebox{5}[13]{$\}$}}	%draw bracket from top to bottom
	%this has been a p.i.t.a. when resizing; there has to be a better way
\put(385,145){\scalebox{2}[2]{\interval{i}}}
\put(10,200){\shortstack{time\\flows\\down}}
\put(5,230){\vector(0,-1){30}}

\end{picture}

 	%generic lattice
%\caption{Generic Lattice}
%\label{fig:genlattice}
%\end{wrapfigure}
%%\end{figure}
%
%\pn Because lattices will be used to define intervals, when an unspecified-further lattice needs a name it is often called \interval{i}.  Interval 
%\(\interval{i}=\langle V_{\interval{i}},E_{\interval{i}}\rangle\)
%%\index{ i@\interval{i}} 
%has a least element at the top called \($start$(\interval{i})\in V_{\interval{i}}\), and an upper bound at the bottom called \($end$(\interval{i})\in V_{\interval{i}}\).  Like trees and conventional current, representations are reflected; least is top (because it's first) and upper most is bottom (because it's last).  To define a notion of ``before" is why all that stuff about irreflexive partial orders, least elements and upper bounds was needed.  
%
%
%\pn Every edge and vertex in the lattice is reachable from the least element; every edge and vertex in the lattice can reach the upper bound.\footnote{$\forall$ means ``for all"}\index{ forall@$\forall$ for all}  
%\(\forall v\in V_{\interval{i}}\minus\ell~\vert~\ell\sqsubset v
%\ \wedge\ \forall v\in V_{\interval{i}}\minus u~\vert~ v\sqsubset u\)
%
%\pn If there is a path between 
%$v_1$ and $v_2$, then $v_{1}\sqsubset v_{2}$, which means $v_1$ occurs \bemph{before} $v_2$.
%If there is no path between $v_1$ and $v_2$, 
%\(v_{1}\nsqsubset v_{2} \wedge v_{2}\nsqsubset v_{1} \rightarrow v_{1}\parallel v_{2}\)
%\index{ concurrent@$\parallel$ concurrent}
%then  $v_1$ and $v_2$ may occur in either order, or concurrently.
%
%\pagebreak 
%\setlength{\unitlength}{0.75pt}
%\renewcommand{\nodetextscalefactor}{0.75} %scales text with unitlength	
%
%%\begin{figure}
%\begin{wrapfigure}{l}{0pt}
%%lattice2.tex
%draw lattices (V1,E1) and (V2,E2)

\setcounter{lwidth}{180}
\setcounter{lheight}{150}

\begin{picture}(\value{lwidth},\value{lheight})(0,0)
%\graphpaper(0,0)(\value{lwidth},\value{lheight})

\setcounter{circlediameter}{15}
\thicklines
\color{blue}

\drawarc{1}{11}{1}{1}{1}{4}
\drawarc{1}{11}{1}{1}{2}{4}
\drawarc{1}{11}{1}{1}{3}{4}
\drawarc{1}{11}{1}{1}{4}{4}

\drawarc{4}{11}{1}{4}{1}{1}
\drawarc{4}{11}{2}{4}{1}{1}
\drawarc{4}{11}{3}{4}{1}{1}
\drawarc{4}{11}{4}{4}{1}{1}

\drawarc{7}{11}{1}{1}{1}{4}
\drawarc{7}{11}{1}{1}{2}{4}
\drawarc{7}{11}{1}{1}{3}{4}
\drawarc{7}{11}{1}{1}{4}{4}

\drawarc{10}{11}{1}{4}{1}{1}
\drawarc{10}{11}{2}{4}{1}{1}
\drawarc{10}{11}{3}{4}{1}{1}
\drawarc{10}{11}{4}{4}{1}{1}

\color{green}
\drawtextnode %{row}{numrows}{node}{numnodesinrow}{text}
	{1}{11}{1}{1}{$\ell_1$}
\drawtextnode{5}{11}{1}{1}{$u_1$}

\drawtextnode{7}{11}{1}{1}{$\ell_2$}
\drawtextnode{11}{11}{1}{1}{$u_2$}

\color{black}
\put(-20,137){\scalebox{\nodetextscalefactor}[\nodetextscalefactor]
	{$\intervalsub{i}{1}=(V_{1},E_{1})$}}
\put(-20,67){\scalebox{\nodetextscalefactor}[\nodetextscalefactor]
	{$\intervalsub{i}{2}=(V_{2},E_{2})$}}

\end{picture}
	%(V1,E1) and (V2,E2)
%\caption{Two Lattices} % $(V_{1},E_{1})$ and $(V_{2},E_{2})$}
%\label{fig:twolattices}
%%\end{figure}
%\end{wrapfigure}
%
%\pn Lattices may combined into new lattices\index{lattice} in three ways:  sequential, concurrent, and insertion.
%Consider two lattices $\intervalsub{i}{1}=\langle V_{1},E_{1}\rangle$ and $\intervalsub{i}{2}=\langle V_{2},E_{2}\rangle$ 
%that have no vertices in common, $V_{1}\cap V_{2} = \emptyset$, 
%least elements $\ell_{1}\in V_{1}$ and $\ell_{2}\in V_{2}$, 
%and upper bounds $u_{1}\in V_{1}$ and $u_{2}\in V_{2}$.
%
%\pn Their \bemph{sequential lattice combination},
%$\intervalsub{i}{1}\curvearrowright\intervalsub{i}{2}$,\index{ seq@$\curvearrowright $ sequential combination} 
%may be performed as follows:  
%substitute $u_1$ for $\ell_2$ in $V_2$ and $E_2$, 
%then form the union of the vertices and edges, $\intervalsub{i}{sc}=\langle V_{1}\cup V_{2},E_{1}\cup E_{2}\rangle$.
%
%\pn Their \bemph{concurrent lattice combination},
%$\intervalsub{i}{1}\downdownarrows\intervalsub{i}{2}$,\index{ con@$\downdownarrows $ concurrent combination} may be performed as follows: 
%substitute $u_1$ for $u_2$, and $\ell_1$ for $\ell_2$ in $V_2$ and $E_2$, 
%then form the union of the vertices and edges, $\intervalsub{i}{cc}=\langle V_{1}\cup V_{2},E_{1}\cup E_{2}\rangle$.
%
%\pn Their \bemph{insertion combination}, $\intervalsub{i}{1}[e=\intervalsub{i}{2}]$,
% may be performed as follows:  
%choose an edge $e\in E_1$, $e=\langle v_{j},v_{k}\rangle$,
%remove it from from $E_1$,
%substitute $v_j$ for $\ell_2$, and $v_k$ for $u_2$ in $V_2$ and $E_2$,
%then form the union of the vertices and edges, $\intervalsub{i}{ic}=\langle V_{1}\cup V_{2},E_{1}\cup E_{2}\rangle$.
%
%
%\begin{figure}
%%\begin{wrapfigure}{r}{0pt}
%%lattice_combinations.tex
%draw lattices for insertion sequential and concurrent combination

\setcounter{lwidth}{600}
\setcounter{lheight}{150}

\begin{picture}(\value{lwidth},\value{lheight})(0,0)
%\graphpaper(0,0)(\value{lwidth},\value{lheight})

\setcounter{circlediameter}{15}
\thicklines

%DRAW LATTICES INSERTION COMBINATION
\color{blue}

\drawarc{1}{6}{1}{7}{1}{16}
\drawarc{1}{6}{1}{7}{2}{16}
\drawarc{1}{6}{1}{7}{3}{16}
\drawarc{1}{6}{1}{7}{4}{16}

\drawarc{2}{6}{2}{16}{3}{40}
\drawarc{2}{6}{2}{16}{4}{40}
\drawarc{2}{6}{2}{16}{5}{40}
\drawarc{2}{6}{2}{16}{6}{40}

\drawarc{4}{6}{3}{40}{2}{16}
\drawarc{4}{6}{4}{40}{2}{16}
\drawarc{4}{6}{5}{40}{2}{16}
\drawarc{4}{6}{6}{40}{2}{16}

\drawarc{5}{6}{1}{16}{1}{7}
\drawarc{5}{6}{2}{16}{1}{7}
\drawarc{5}{6}{3}{16}{1}{7}
\drawarc{5}{6}{4}{16}{1}{7}

\color{green}
\drawtextnode %{row}{numrows}{node}{numnodesinrow}{text}
	{1}{6}{1}{7}{$\ell_1$}

\drawtextnode{2}{6}{2}{16}{$\ell_2$}
\drawtextnode{5}{6}{2}{16}{$u_2$}

\drawtextnode{6}{6}{1}{7}{$u_1$}

\color{black}
\put(50,0){\scalebox{\nodetextscalefactor}[\nodetextscalefactor]
	{$\intervalsub{i}{1}[e=\intervalsub{i}{2}]$}}

%\color{black}
%\put(-20,137){\scalebox{\nodetextscalefactor}[\nodetextscalefactor]
%	{$\intervalsub{i}{1}=(V_{1},E_{1})$}}
%\put(-20,67){\scalebox{\nodetextscalefactor}[\nodetextscalefactor]
%	{$\intervalsub{i}{2}=(V_{2},E_{2})$}}

%SEQUENTIAL COMBINATION
\color{blue}
\drawarc{1}{11}{3}{7}{5}{16}
\drawarc{1}{11}{3}{7}{6}{16}
\drawarc{1}{11}{3}{7}{7}{16}
\drawarc{1}{11}{3}{7}{8}{16}

\drawarc{4}{11}{5}{16}{3}{7}
\drawarc{4}{11}{6}{16}{3}{7}
\drawarc{4}{11}{7}{16}{3}{7}
\drawarc{4}{11}{8}{16}{3}{7}

\drawarc{5}{11}{3}{7}{5}{16}
\drawarc{5}{11}{3}{7}{6}{16}
\drawarc{5}{11}{3}{7}{7}{16}
\drawarc{5}{11}{3}{7}{8}{16}

\drawarc{9}{11}{5}{16}{3}{7}
\drawarc{9}{11}{6}{16}{3}{7}
\drawarc{9}{11}{7}{16}{3}{7}
\drawarc{9}{11}{8}{16}{3}{7}

\color{green}
\drawtextnode %{row}{numrows}{node}{numnodesinrow}{text}
	{1}{11}{3}{7}{$\ell_1$}
\drawtextnode{5}{11}{3}{7}{$u_1=\ell_2$}
%\drawnode{5}{11}{3}{7}

\drawtextnode{10}{11}{3}{7}{$u_2$}

\color{black}
\put(210,0){\scalebox{\nodetextscalefactor}[\nodetextscalefactor]
	{$\intervalsub{i}{1}\curvearrowright\intervalsub{i}{2}$}}
	
%CONCURRENT COMBINATION
\color{blue}
\drawarc{1}{11}{6}{7}{9}{16}
\drawarc{1}{11}{6}{7}{10}{16}
\drawarc{1}{11}{6}{7}{11}{16}
\drawarc{1}{11}{6}{7}{12}{16}
\drawarc{1}{11}{6}{7}{13}{16}
\drawarc{1}{11}{6}{7}{14}{16}
\drawarc{1}{11}{6}{7}{15}{16}
\drawarc{1}{11}{6}{7}{16}{16}

\drawarc{4}{11}{9}{16}{6}{7}
\drawarc{4}{11}{10}{16}{6}{7}
\drawarc{4}{11}{11}{16}{6}{7}
\drawarc{4}{11}{12}{16}{6}{7}
\drawarc{4}{11}{13}{16}{6}{7}
\drawarc{4}{11}{14}{16}{6}{7}
\drawarc{4}{11}{15}{16}{6}{7}
\drawarc{4}{11}{16}{16}{6}{7}


\color{green}
\drawtextnode %{row}{numrows}{node}{numnodesinrow}{text}
	{1}{11}{6}{7}{$\ell_1=\ell_2$}
\drawtextnode{5}{11}{6}{7}{$u_1=u_2$}
%\drawnode{5}{11}{3}{7}


\color{black}
\put(430,60){\scalebox{\nodetextscalefactor}[\nodetextscalefactor]
	{$\intervalsub{i}{1} \downdownarrows\intervalsub{i}{2}$}}



\end{picture}
	%(V1,E1) and (V2,E2)
%\caption{Lattice Combinations} % $(V_{1},E_{1})$ and $(V_{2},E_{2})$}
%\label{fig:latticecombinations}
%\end{figure}
%%\end{wrapfigure}

%MEANING
%\section{Meaning}\label{sec:informal_meaning}
%
%\pn The \bemph{meaning} of \lang~ language constructs is defined by giving an \emph{interpretation} 
%within a \emph{context} for a \emph{subject} by a \emph{model}: \\
%\(\llbracket\mathfrak{M}_{context}\models~$subject$\rrbracket\equiv~$interpretation$\) \\
%where\index{ meaning@$\mathfrak{M}$ meaning}  %\index{ meaningx@$\vert[ \vert]$} 
%\begin{description}
%%\begin{indent}
%	\item[$\mathfrak{M}$] is the model which determines meaning of symbols
%	\item{$\models$} ``models''
%	\item[context] if given, is usually an interval lattice, a state in a lattice, or a moment in time
%	\item[subject] is some construct in \lang~
%	\item[interpretation] is the defining formula for that subject, in that context
%%\end{indent}
%\end{description}

%TIME
\section{Time}\label{sec:informal_timemodel}

\pn It is important to distinguish \bemph{model time} from \bemph{real time}.  
Model time is a mathematical abstraction useful for defining what a system is supposed to do.
Real time is where the actual systems will operate.
Model time is the same everywhere in the system, and is non-negative and real, $t\in\mathbb{R}\wedge 0 \le t \le now$.
Real time is different everywhere; synchronous temporal domains are limited in volume by the speed of light.
Great care must be used for information that flows across temporal domain boundaries.
Defining such temporal domains in AADL using precise, model time is means to make them nicely work together in real time when integrated into an operational system.

\pn Periods discretize time exactly in \lang~ by choosing a subset of $\mathbb{R}$,
\[P_{d}=\{p_{j}~\vert~ p_{j}=dj~$for all$~j\in\mathbb{N}_{0}\}\]
where $d$ is the period's duration, and $dj$ is multiplication of $d$ by $j$.

\pn Defining the system's \bemph{hyperperiod} is naturally the product of all of the different 
durations:\footnote{In cases where every system-level clock is a multiple of the same reference clock, then least-common multiple of different durations can suffice.}
\[h\equiv\prod_{d\in D} d\]
where $D$ is the set of all different durations in the system.\footnote{$\prod$ means ``product of", usually defined over all of the numbers in a given set}\index{ productof@$\prod$ product of}  

\pn The present instant is called \bemph{now}.

%\pn Many entities in \lang~ have sensible time of occurrence, %$\mathfrak{T}$, 
%$T$, such as events.
%\[T\llbracket $e$\rrbracket~\equiv~t~ ~\vert~~t\in\mathbb{R}\wedge 0 \le t \le now \wedge ~$e occurs at$~t\]

\pn   Durations are continuous sets of non-negative real numbers.  Commas denote open ranges that do not include the upper and/or lower bound: ``.."=closed, includes both endpoints;
``,,"=open both, neither endpoint included;
``,."=open left, lower bound not included;
``.,"=open right, upper bound not included.
\index{ xdd@\lstinline+..+ closed interval}
\index{ xcd@\lstinline+,.+ open left}
\index{ xdc@\lstinline+.,+ open right}
\index{ xcc@\lstinline+,,+ open interval}

\(
\begin{array}{l}
\{l..u\}~\equiv~ \{~m~~\vert~~m\in\mathbb{R}\wedge m\ge l\wedge m\le u~\} \\
\{l,,u\}~\equiv~ \{~m~~\vert~~m\in\mathbb{R}\wedge m>l\wedge m<u~\} \\
\{l,.u\}~\equiv~ \{~m~~\vert~~m\in\mathbb{R}\wedge m<l\wedge m\le u~\} \\
\{l.,u\}~\equiv~ \{~m~~\vert~~m\in\mathbb{R}\wedge m\ge l\wedge m<u~\}
\end{array}
\)

%\pn Frequently, time will be used to define the context of meaning.  Subscripted time as context notation
%\(\mdx{X}{t}~\)
%is used to define the meaning for whatever $X$ is, at a given time $t$.  This notation is used in sections
%\ref{sec:timedpredicate} and \ref{sec:timedexpression} to define temporal meaning for Assertions.
%\footnote{Assertions would otherwise be simply first-order predicates.}


\section{Values}\label{sec:informal_values}

\pn A \bemph{value} is a mathematical object.  A \bemph{type} is a set of values. % (see D\ref{chapter:type} Type).  
Values used in \lang~ are the same as the AADL Data Modeling Annex and AADL property values (which have records, but not arrays).


\pn Usually, values are singular:  quantities in $\mathbb{Z}$ (integers), or $\mathbb{R}$ (real numbers); boolean in $\mathbb{B}$; a character (enclosed in apostrophes); a string (enclosed in quotation marks); or an enumeration literal (sequence of alphanumeric characters starting with a letter).  

\pn More complex values are constructed from sets of pairs.  Records are sets of pairs in which the first element is a record field identifier, and the second is the value of that field.  Arrays are sets of pairs in which the first element is an integer index, and the second is the value of that index.  Record and array values are functions in which the first element of the pair is unique.  Array values generally constrain the first element of pair to me a range (consecutive) of integers, and  the second element of pairs to the same type.  Of course, array element and record field values can be themselves be arrays or records, making arbitrarily complex values.

\pn The empty set, $\emptyset$, represents the absence of a value at a given time.  The absence of value is also called \bemph{null}.

%\JP{
%\pn The \bemph{clock operator} $\hat{}$ determines when something has value.  At time $t$,
%
%\(\mathfrak{M}_{t}\llbracket\hat{p}\rrbracket\equiv~
% \begin{array}{l}
%false$ when $\mathfrak{M}_{t}\llbracket$p$\rrbracket=\bot \\
%true $ otherwise$
% \end{array}
%\)
%}

%STATES
%\section{States}\label{sec:informal_states}
%
%\pn \lang~ uses two kinds of states:
%\begin{description}
%\item[lattice states] variable-value bindings during actions
%\item[machine states] source and destination of transitions
%\end{description}
%
%\subsection{Lattice States}\label{subsec:latticestates}
%
%\pn A \bemph{lattice state} is a set of pairs of variable names with values, with perhaps a time of the moment of occurrence.  
%
%%\indent
%\[L=(\{s_{1},s_{2},\ldots,s_{m}\},t_S)~~s_{k}=\langle $n$_{k},$v$_{k}\rangle~~t_S \in \mathbb{R}\wedge t_S\ge0\]% \\
%%or
%%\[S=(\{($n$_{1},$v$_{1}),($n$_{2},$v$_{2}),\ldots,($n$_{m},$v$_{m})\},t_S)\]
%
%Two states are equal if-and-only-if they have 
%the same variables and those variables have the same values, but not necessarily the same time of occurrence. 
% For states $V$ and $U$,
%\[ V=(\{v_{1},v_{2},\ldots,v_{m}\},t_V)~$and$~U=(\{u_{1},u_{2},\ldots,u_{m}\},t_U)  \]
%\[V=U\equiv \{v_{1},v_{2},\ldots,v_{m}\}=\{u_{1},u_{2},\ldots,u_{m}\}\]
%
%Execution lattices satisfying temporal logic formulas have states as vertices (nodes).
%%put formula for equality of states here?
%
%
%\subsection{Behavior States}\label{subsec:behaviorstates}
%
%\pn A \bemph{behavior state} is declared in the \lstk{states} section of thread behaviors (section \ref{sec:behaviorstates}), and may be used as sources or destinations of transitions.
%
%\[Q=(s,L)\]
%
%Where $s$ is a behavior state label, and $L$ is a lattice state defining values of variables and time of occurrence. 

\section{Arithmetic}\label{sec:informal_Arithmetic}

\pn Axiomatic definitions of arithmetic have been of interest to mathematicians for centuries, but only rigorously derived from logic and set theory in the early 20th century.  For \lang~, $\mathbb{Z}$ (integers) and $\mathbb{R}$, use definitions and derivations in Metamath \texttt{set.mm} database (see \ref{sec:mmnumbers}).

\section{Logic}\label{sec:informal_logic}

\pn A \bemph{logic} is formal mathematical system for reasoning about a domain of interest.  A logic consists of rules defined in terms of
\begin{description}
\item[symbols] a set of graphical characters
\item[formulas] a set of sequences of symbols
\item[axioms] a set of distinguished formulas known to be true 
\item[rules] a set of inferences to prove additional formulas from axioms, given formulas, and previously proved formulas
\end{description}
Not all sequences of symbols are formulas.  
Formulas are \bemph{well-formed} sequences of symbols.  Formulas must be grammatically-correct to have meaning. A logic usually defines which sequences of symbols are formulas with a grammar.

To \bemph{satisfy} a formula means choosing values for its symbols such that the formula is true.  A formula for which no choice of values for symbols is true is \bemph{unsatisfiable}.  A formula always true regardless of chosen values for symbols is \bemph{tautology}.

%The following formulas are assumed as axiomatic (eg. tautologies), with $b$, $c$, and $d$ representing boolean-valued predcates, $r$ being a bounded range, and $j$ being an element in that range:\footnote{as in U.S. Pat. No. 5,867,649, col. 40, lines 40-70}
%
%\begin{law}[\bemph{Complement}]\label{lawofcomplement} 
%\(b \equiv \lnot(\lnot b) \)
%\end{law}
%\index{ complement@$\neg$ complement}
%\begin{law}[\bemph{Excluded Middle}]\label{lawofexcludedmiddle} 
%\(b \lor \lnot b\)
%\end{law}
%\begin{law}[\bemph{Contradiction}]  \label{lawofcontradiction}
%\(\lnot(b \land \lnot b)\)
%\end{law}
%\begin{law}[\bemph{Implication}]\label{lawofimplication}  
%\(a\rightarrow b \equiv \lnot a \lor b \)
%\end{law}
%\index{ implication@$\rightarrow$ implication}
%\begin{law}[\bemph{Equality}]  
%\(a=b \equiv a\rightarrow b \lor b\rightarrow a\)\label{lawofequality}
%\end{law}
%\begin{law}[\bemph{Disjunction}]\label{lawofdisjunction1} 
%  \(c\lor c \equiv c\) 
%\end{law}
%\begin{law}[\bemph{Disjunction}]\label{lawofdisjunction2} 
%  \(c\lor true \equiv true\) 
%\end{law}
%\begin{law}[\bemph{Disjunction}]\label{lawofdisjunction3} 
%  \(c\lor false \equiv c\) 
%\end{law}
%\begin{law}[\bemph{Disjunction}]\label{lawofdisjunction4} 
%  \(c\lor (c\land b) \equiv c\) 
%\end{law}
%\begin{law}[\bemph{Disjunction}]\label{lawofdisjunction5} 
%  \(b\lor c \equiv c\lor b \)  
%\end{law}
%\begin{law}[\bemph{Disjunction}]\label{lawofdisjunction6} 
%  \(b \rightarrow (b\lor c)  \)
%\end{law}
%\index{ disjunction@$\vee$ disjunction}
%\begin{law}[\bemph{Conjunction}]\label{lawofconjunction1}  
%  \(b\land b \equiv b\) 
%\end{law}
%\begin{law}[\bemph{Conjunction}]\label{lawofconjunction2}  
%  \(b\land true \equiv b\) 
%\end{law}
%\begin{law}[\bemph{Conjunction}]\label{lawofconjunction3}  
%  \(b\land false \equiv false\) 
%\end{law}
%\begin{law}[\bemph{Conjunction}]\label{lawofconjunction4}  
%  \(b\land (c\lor b) \equiv b\) 
%\end{law}
%\begin{law}[\bemph{Conjunction}]\label{lawofconjunction5}  
%  \(b\land c \equiv c\land b \)  
%\end{law}
%\begin{law}[\bemph{Conjunction}]\label{lawofconjunction6}  
%  \((b\land c) \rightarrow b  \)
%\end{law}
%\index{ conjunction@$\wedge$ conjunction}
%\begin{law}[\bemph{Distribution}]\label{lawofdistribution1}  
%\(b \lor (c\land d) \equiv (b\lor c)\land(b\lor d)\) 
%\end{law}
%\begin{law}[\bemph{Distribution}]\label{lawofdistribution2}  
%\(b \land (c \lor d) \equiv (b \land c) \lor(b \land d)\)
%\end{law}
%\begin{law}[\bemph{Universal Quantification}]\label{lawofuniversalquantification}  
%  \(\forall j \in r~\vert~(b \land c)~\equiv~(\forall j \in r~\vert~b)\land(\forall j \in r ~\vert~ c)  \)
%\end{law}
%\index{ forall@$\forall$ for all}
%\begin{law}[\bemph{Existential Quantification}]\label{lawofexistentialquantification}  
%  \(\exists j \in r~\vert~(b \lor c)~\equiv~(\exists j \in r~\vert~b) \lor(\exists j \in r ~\vert~ c)  \)
%\end{law}
%\index{ bar@$~\vert~$ such that}
%\index{ exists@$\exists$ exists}

%\section{Computation $\equiv$ Satisfaction}\label{sec:informal_satisfaction}
%
%\pn \lang~ defines computation as satisfaction of interval temporal logic formulas by lattices of states:
%
%\(\mdi{w}\)~\ssc{construct an interval \interval{i} such that the formula $w$ is true} %\index{ true@$\top$ true}
%
%\pn Each \lang~ program may be satisfied by a huge number of different lattices, all of which arrive at the same 
%result.\footnote{Every satisfying lattice will have equal states for their least elements (start) and upper bounds (end).}  The set of satisfying lattices is so large, it is effectively countably infinite.  However, it suffices to consider the canonical member of the set of satisfying lattices--the shortest and bushiest lattice.  Maximizing opportunities for concurrent execution is paramount for supercomputing, but embedded systems with multi-core systems-on-chip may benefit from rich opportunities for concurrent execution.

%\pn For just the Action annex sublanguage (chapter \ref{chapter:action}), the states\index{state} defined in section \ref{sec:states} suffice and need no more reference in time than its position in the lattice.  For satisfying lattices of states for the \lang~ annex sublanguage need time-stamps.  Therefore, the set of variable-value pairs comprising a state is augmented with a real-valued time-stamp.
%\[L=(\{s_{1},s_{2},\ldots,s_{m}\}, t_{s})~~s_{k}=\langle $n$_{k},$v$_{k}\rangle\]
%%S=(\{($n$_{1},$v$_{1}),($n$_{2},$v$_{2}),\ldots,($n$_{m},$v$_{m})\}, t_{s}\]
%where $t_{s}$ the time lattice state $L$ is created.  Lattice state $L$ says nothing about the values of variables at any time other than $t_{s}$.
%Other lattice states could have occurred infinitesimally earlier or later.  Usually, only the time-stamps of least elements and upper bounds of lattices matter, and will be ignored when they don't.

%\section{Proofs}\label{sec:informal_proofs}
%
%\pn A mathematical \bemph{proof} is a sequence of formulas, each of which is an axiom of the logic\footnote{see D\ref{sec:logic}} used, a given assumption, or a theorem derived by an inference rule from axioms, assumptions, or other theorems earlier in the sequence.  The last formula is usually the desired truth to be established such as the program meets its specification and therefore is correct.
%
%\pn Proofs for \lang~ are sequences of triples\footnote{often called ``Hoare" triples for Tony Hoare who originatied them.} of the form:  \lstk+<<P>>+~\lstk+S <<Q>>+, where \lstk+<<P>>+ and \lstk+<<Q>>+ are Assertions, and \lstk+S+ is behavior actions.%\footnote{\texttt{behavior\_actions} is the root production of the middle term of every triple}.  
%\lstk+<<P>>+ is the \bemph{precondition} which must be true before performing the behavior actions \lstk+S+.  \lstk+<<Q>>+ is the \bemph{postcondition} which must be true after performing the behavior actions \lstk+S+.  
%
%\pn When there is no action, the triple becomes \lstk+<<P>>+~\lstk+-> <<Q>>+.  Does \lstk+P+ imply \lstk+Q+?  Equivalently, \(\inference{$\lstk+<<P>>+$}
%%__________________________?
%              {$\lstk+<<Q>>+$}\)? asks whether Q can be derived from P.
%
%\pn Each triple in the proof must have a \bemph{reason}:  an axiom (or given), or an inference rule applied to triples earlier in the sequence.  The last triple in the sequence is a mathematical statement that a program meets its specification.
%
%%\pn When life was easy, back in the days when computing was the process of transforming data from the state you got it to the state you want it, \lstk+<<P>> S <<Q>>+ was first-order predicates bracketing sequential, imperative, and occasionally non-deterministic programs.  Back then, subprogram proofs had a single proof obligation.  
%
%\pn AADL subprogram components that define behavior with the Action language need a single proof obligation.  AADL thread components on the other hand, have \emph{dozens} of proof obligations, and the Assertions are first-order predicates augmented expressions of time to which the predicates apply.
%
%\section{Induction}\label{sec:informal_Induction}
%
%\pn Proofs often rely on \bemph{induction}.  Usually induction involves proving some property $P$, 
%that has a natural number parameter, say $n$.
%To prove that P holds for all $n\in\mathbb{N}_0$, it suffices to proceed by induction on $n$,
%organizing the proof as follows:
%\begin{itemize}
%\item induction basis:  prove that $P$ holds for $n=0$.
%\item induction step:  prove that P holds for $n+1$ from the induction hypothesis that $P$ holds for $n$.
%\end{itemize}
%The induction principle for natural numbers is based on the fact that natural numbers can be constructed beginning with the number $0$ and repeatedly adding $1$.  
%
%\pn Proofs of loops use induction over a loop invariant.  Essentially, \lang~ behavior is induction over time.  
%%For the VVI example (D\ref{chapter:vvi}), this induction is over cardiac cycles, which may vary in length, for the duration of operation (6+ years).  
%Induction is very powerful, used often, is tricky, and deserves scrutiny.

